% ID: SEL-D0029-OB-E03
% ESTADO: validada. S030, cotejada con examen y criterios.
Los tres grupos suman 250. Se toman como categorías disjuntas,
según el reparto del enunciado. Hay $120(0.75)=90$ aficionados
al ciclismo en fútbol, $60(0.65)=39$ en tenis y $70(0.60)=42$
en baloncesto.

\textbf{a)} De los 250 deportistas, $90+39+42=171$ son aficionados
al ciclismo. La probabilidad es $\boxed{171/250=0.684}$.

\textbf{b)} Al condicionar, el total pasa a ser 171. De ellos 39
juegan al tenis, de modo que la probabilidad es
$\boxed{39/171=13/57\approx0.2281}$.
